\section{Introduction}

Adult zebrafish telencephalon and optic tectum do very different work. The telencephalon participates in social behavior, learning, memory, and other forebrain functions \citep{stednitz2018social,reemst2023memory,pandey2023telencephalon}. The optic tectum, by contrast, is a visuomotor center that transforms sensory input into orienting and avoidance behavior \citep{roeser2003tectum,orger2016visuomotor}. If that functional division is written into the protein state of each region, it should be visible most clearly at the proteoform level, where sequence variants and post-translational modifications are resolved rather than folded into protein totals.

The LCM-CZE-MS/MS study of \citet{lubeckyj2022lcmcze} offers an unusual opportunity to ask that question directly. The article reported hundreds of proteoforms from adult zebrafish telencephalon and optic tectum and linked those profiles to regional function. Just as important, the study deposited two region-specific spreadsheets, \texttt{Proteoform\_ID\_Tel2.xlsx} and \texttt{proteoform\_ID\_Teo2.xlsx}, in the PRIDE proteomics repository. Those files expose the overlap structure of the dataset in a way that the narrative of the paper alone cannot. They allow direct comparison of exact accession+proteoform identities, duplicate incidence, protein-level collapse, and biologically interpretable marker families from the deposited evidence itself.

Read in that form, the dataset makes a sharper statement than the original article drew explicitly. Telencephalon and optic tectum share very few proteoforms, and the proteoforms that distinguish them fall along a biologically intelligible axis: neuropeptide and synaptic entries on the telencephalic side, reticulon- and cytoskeleton-linked entries on the tectal side. The apparent excess of N-terminal acetylation in telencephalon is much less secure and weakens once uneven recovery is taken into account. The analysis presented here was designed around that distinction. It asks how strong the regional separation remains under exact matching and conservative sensitivity analysis, whether the separating proteoforms align with known regional biology, and where the evidence becomes too thin to support a firm claim.

This question matters for more than one zebrafish dataset. Prior functional studies and atlases establish the broad biology of the two regions \citep{stednitz2018social,pandey2023telencephalon,roeser2003tectum,orger2016visuomotor}, and recent top-down work shows that proteoform-level structure in zebrafish brain remains informative well beyond a single pilot study \citep{askarani2025aging}. At the same time, region-resolved datasets are often discussed at the level of appealing examples rather than explicit overlap structure. The present study shows that the deposited proteoform evidence is already strong enough to support a definite statement about regional molecular organization in the adult zebrafish brain.
